% ================================================================= 14
\section{Limitations}\label{sec:limits}
\begin{itemize}
\item \textbf{Sample size.} 82, 80, 44, 17 and 10 questions show large differences, not small ones; per-type results and
the Zep comparison are descriptive.
\item \textbf{Thresholds.} The 0.9 floor follows the calibration analysis on LongMemEval decisions, so BEAM is the
held-out test. Thresholds are set, not fitted on human labels, and agreement with Sonnet is not ground truth. A
labelling tool and a fitting script are published; the 300 human labels are not yet collected.
\item \textbf{Models.} Memory systems extract with Haiku 5.5 throughout; OrgProbe used Haiku 4.5 for parity with its
mem0 run. Numbers compare within a table, not across tables.
\item \textbf{Ingest cost.} \mx{} stores more facts than mem0 and its ingest costs about the same (Section~\ref{sec:lme});
the saving is at every recall. HopProbe and OrgProbe are author-written.
\item \textbf{A hosted, closed decision model.} A Jev update can move the probability scale under fixed thresholds;
\mx{} pins \code{jev-1.14} and stores the question version with every decision.
\item \textbf{Conflict of interest.} \mx{} is the author's company's product. Harnesses, fixtures and result files are
published.
\end{itemize}

% ================================================================= 15
\section{Conclusion}
\mx{} puts a calibrated judge at the centre of agent memory: a generative model writes facts, Jev decides what happens
to them, and code applies asymmetric thresholds and logs every answer. Scoped corrections, a five-way verdict with a
0.9 floor on destructive outcomes, asymmetric entity priors and joint grading of linked facts give a memory that is
small at recall time and more accurate than mem0 at the same ingest cost: 2--9 times fewer tokens per answer, 89\%
from full memory on LongMemEval, the best score on BEAM, and a walk that answers two-hop questions search cannot. Next
steps are a subject name for ``the user'' so personal facts link to the person entity, and thresholds fitted on the
human-labelled decisions the published tool collects.


\section*{Availability}
Code, harnesses (\code{eval/longmemeval}, \code{eval/beam}, \code{eval/orgprobe}, \code{eval/hopprobe},
\code{eval/calibration}), result files and the figure script
are in the \mx{} repository (\code{docs/paper/tech-report}).
